基于对0-Intro.v2_polished_unified_expanded.tex文件的分析，以下是我诊断出的需要扩写1-1.5倍的段落及其扩写版本：

【原文1】
这种由"困难"走向"显而易见"的转变并非偶然，而是研究范式成熟的结果。当我们抓住了问题的本质，并匹配到恰当的数学工具与计算方法，看似不可逾越的障碍就会自然消解。

【扩写版本1】
这种由"困难"走向"显而易见"的转变绝非偶然现象，而是研究范式逐步走向成熟的必然结果。为什么这样说呢？关键在于我们对问题本质的理解在不断深化。当我们真正抓住了问题的核心本质，并为之匹配到恰当的数学工具与计算方法时，那些原本看似不可逾越的障碍就会像冰雪消融般自然化解。

这就像学习骑自行车的过程。初学者觉得平衡与方向控制极其困难，每一个动作都需要刻意思考。但一旦掌握了"重心转移"和"身体协调"的本质要领，骑车就变成了近乎本能的行为。同样，在AI发展中，许多技术突破都源于这种"本质洞察"——当我们找到了描述和解决问题的正确视角，原本复杂的问题就变得清晰可解。

【原文2】
线性代数让我们得以处理高维数据，微积分使得梯度下降成为可能，概率论为不确定性建模奠定了基础，而信息论则揭示了学习与压缩的内在联系。回望人工智能的发展，每一次重大突破几乎都伴随着数学思想的革新或重生。

【扩写版本2】
线性代数为我们打开了处理高维数据的大门，让计算机能够理解和操作复杂的向量空间；微积分使得梯度下降等优化算法成为可能，为机器学习提供了寻找最优解的数学工具；概率论为处理现实世界的不确定性建模奠定了坚实基础，让AI系统能够在信息不完整的条件下做出合理推断；而信息论则深刻揭示了学习与压缩之间的内在联系，帮助我们理解模型泛化的本质。

如果我们仔细回望人工智能的发展历程，会发现一个有趣的规律：每一次重大技术突破几乎都伴随着数学思想的革新或重生。这就像盖房子需要不同的建筑材料，不同的数学工具为AI的发展提供了不同的支撑能力。没有线性代数，我们就无法有效处理图像、文本等高维数据；没有微积分，梯度优化算法就无从谈起；没有概率论，我们就无法处理现实世界的不确定性；没有信息论，我们就无法理解学习的本质。

【原文3】
深度学习的成功，本质上是这些数学思想的再整合。神经网络的反向传播算法体现了微积分的链式法则，卷积神经网络继承了信号处理的卷积理论，而注意力机制则借鉴了概率分布与加权平均的思维。

【扩写版本3】
深度学习的惊人成功，本质上可以看作是这些经典数学思想在现代计算环境下的重新整合与升华。这个现象为什么值得深入思考？因为它揭示了不同数学分支之间的内在统一性。

具体来说，神经网络的反向传播算法完美体现了微积分中的链式法则——通过层层传递误差信号，实现了复杂函数的高效优化；卷积神经网络巧妙继承了信号处理中的卷积理论，将局部连接和权值共享的思想发挥到极致；而近年来大放异彩的注意力机制，则借鉴了概率分布与加权平均的思维，让模型能够动态地关注输入信息中的关键部分。

这种融合不是简单的拼接，而是深层次的数学统一。就像一位技艺精湛的厨师，将不同食材的风味完美融合，创造出全新的味觉体验。深度学习就是AI领域的"美食烹饪"——它将不同的数学思想和谐地结合在一起，产生了远超各部分简单叠加的强大能力。

【原文4】
最初阶段的人工智能以符号与规则推理为基础。二十世纪中叶的第一代研究者相信，只要能把人类的知识形式化，机器就能像人类一样思考。早期的专家系统与定理证明程序正是这一理念的体现——它们擅长"演绎"，却难以处理模糊性与经验的不确定。

【扩写版本4】
人工智能最初的发展阶段建立在符号与规则推理的坚实基础之上。在二十世纪中叶，第一代AI研究者们怀抱着一个简单而有力的信念：只要能够把人类的专家知识成功地形式化为符号规则，机器就能够像人类一样进行逻辑思考。这种思想为什么在当时如此吸引人？因为它提供了一条清晰的路径来模拟人类智能。

早期的专家系统与定理证明程序正是这一理念的典型体现。这些系统在特定领域展现出了令人印象深刻的能力——它们能够严格地遵循逻辑规则进行推理，在数学定理证明、医疗诊断等任务中取得了不错的效果。然而，这些系统也暴露出了明显的局限性：它们虽然擅长"演绎推理"，却难以处理现实世界中普遍存在的模糊性、矛盾和经验的不确定性。

就像一个只会背诵规则的学生，面对没有明确规则的现实问题时就会束手无策。第一代AI系统就像是这样的学生——在规则明确的世界里表现出色，但在需要灵活应对复杂现实的场合就显得力不从心。这种局限性也促使研究者们开始探索新的智能范式。
